%\addbibresource{/home/jorgsk/Dropbox/phdproject/bibtex/jorgsk.bib}
\subsection{Model of initial transcription}
Our model describes the process of initial transcription by considering the
following  reactions  i) NAC ii) backtracking iii) unscrunching and abortive
RNA release (UAR), and iv) promoter escape (Figure
\ref{fig:model_and_rates}A). Here, we consider backtracking  as the first
unscrunching step, also known as backstepping. The rate constants for these
reactions are estimated using our experimental data.

We assume that the probability to backtrack at any given template position is
equal to the abortive probability (AP) at that position. We used the AP values
that were calculated from experiments by Hsu et al. \cite{hsu_initial_2006}.
For a given freely sampled (WHATS THAT? FREELY SAMPLED RATE?) rate constant of
the NAC, we calculated the rate constant of backtracking as follows:
\begin{equation}
  b_i = \frac{\text{NAC}_i\cdot\text{AP}_i}{1-\text{AP}_i}.
  \label{eq:backtrackingcalc}
\end{equation}
with the subscript $i$ denotes the position of the RNAP. This equation
represents a relationship between the rate of NAC and backtracking (I DISSAGRE
    THAT IT IS COMPETITION, THERE IS NO COMPETITION HERE, NAC [AP/(1-AP)] THE
    RATE WILL BE ZERO WHEN AP=0, AND EXPLODE TO INFINITY WHEN AP APPROACHES 1.
    THE PROBLEM IS THAT SOMEONE WILL COMMENT (I WOULD!) THAT IF AP IS LOW, AND
    NAC IS HIGH, THEN YOU GET THE OPPOSITE EFFECT THAN DESIGRED, BECAUSE HIGH
    NAC SHOULD PREVENT BACKTRACKING, WHEREAS HERE ITS ENCOURGES IT (BY HIGH
    RATE OF b). IT WILL ONLY WORK IF AP AND NAC ARE HIGH OR NAC LOW AND AP IS
WHATEVER.) Figure~\ref{fig:model_and_rates}B shows the initial transcription
kinetic reaction diagram with the rate constants to be estimated. 

In this work we apply the model using APs obtained both for the absence of
GreB (-GreB) and in the presence of GreB (+GreB). When GreB is present,
backtracked complexes may be rescued by GreB stimulated cleavage of the
unaligned 3\ppp end of the transcript. Therefore, APs obtained in the presence
of GreB represent the probability to both backtrack and to avoid rescue by
GreB until abortive RNA release. We assume that GreB-mediated cleavage and
subsequent NACs are rapid steps compared to unscrunching, using the APs
obtained in the presence of GreB as effective backtracking probabilities. All
AP values are obtained from Hsu et al. \cite{hsu_initial_2006}.

\subsection{Implementation and parameter estimation}
The data used for parameter estimation is the distribution of time spent in
abortive cycling \cite{revyakin_abortive_2006}. This data was obtained from
100 individual initial transcription events. For such low copy-number
experiments, there is an inherent stochastic component in the experimental
outcome that derives from the randomness of molecular motion. Indeed, the
rates of transcription elongation obtained by single-molecule experiments show
large standard deviations \cite{adelman_single_2002,
tolic-norrelykke_diversity_2004}. We therefore simulate initial transcription
using stochastic simulations. XXX Specifically, we use the direct method for the
stochastic simulation algorithm \cite{gillespie_exact_1977} as implemented in
the StocPY software \cite{maarleveld_stochpy:_2013}. DETTE MÅ FORKLARES BEDRE,
DU KAN IKKE VÆRE MER DIFFUS ENN DETTE. 1) JEG ANTAR AT MANGE REVIEWER SOM KAN
MODELLERING HAR IKKE HØRT OM STOCHASTIC SIMULATION ALGORITHM STOCPY. KUNNE DU
GI DETALJER MATEMATISK? FOR EKSEMPEL: TO INCORPORATE STOCHASTICITY INTO
INITIATION PROCESS, WE USED (OR THE MOTION IS DESCRIBED BY A RANDOM
VARIABLE..) RANDOM PROBABILITY $(X \sim \mathcal{N} (0,1) )$ TO DESCRIBE THE
MOTION INITIATION.... JEG KAN SKRIVE DET, MEN JEG FORTSATT IKKE VET HVORDAN DU
IMPLEMENTERTE DETTE. 

We estimate rate constants by first randomly sample 100000 triplets of the
rate constants for the NAC, promoter escape, and unscrunching and abortive RNA
release. We then use the value for the NAC to calculate rate constants for
backtracking through equation \eqref{eq:backtrackingcalc}. Having thus obtained
all four rate constants, we simulate 100 initial transcription events to find the
predicted distribution of time spent in abortive cycling specific to these
rate constants. This distribution is then compared to the measured values from
Revaykin et al. \cite{revyakin_abortive_2006} using the root of mean square
error (RMSE). The full procedure is illustrated in Figure
S\ref{fig:parameter_estimation_scheme}. %fixed the numbering to S1. Se
comments on S1.  KOMMENTAR S1: 
VI KOMMER TIL Å FÅ KOMMENTAR HER, FORDI DET 'VANLIGE' PROSEDYREN ER LEAST
SQUARE ELLER MINIMUM LIKELIHOOD. JEG MENER AT I FØLGE FIGURE S1, HAR DU INGEN
OPPDATERING MELLOM RMSE, OG EN ALGORITHM SOM SKAL MINIMALISERE DEN.... ER DET
SLIK AT DU BARE VELGER 100000000 RANDOM VERDIER TIL PARAMETRENE, DA BEREGNER
10000000 RMSE, OG VELGER DEN SOM GAV DEG MINSTE VERDIEN ET AV DISSE 10000000
SIMULAREINGER? DET FINNES BEDRE OG EN STANDARD MÅTE Å GJØRE DEN PÅ:
OPTIMALISERINGS (MINIMALISERING AV RMSE) PROBLEM.  

JEG OM SKREVET PROSEDYREN: We estimated the rate four constants by using the
following procedure (Figure S1): First, we assigned random values to the three
rate constants NAC ($k_n$), promoter escape ($k_e$) and the uncrunching and
abortive RNA release ($k_u$) within certain fixed boundaries (see below).
These values were chosen independently from a uniform distribution ($\sim \cal
U (0,1)$) IS IT UNIFORMLY? OR NORMAL?. Secondly, we calculated the rate
constant of backtracking ($k_b$) using equations \ref{eq:backtrackingcalc}. We
then simulated 100 initial transcription events HER MÅ VI FORKLARE EVENTS!
BEREGNING I EN TILFELDIG SEKVENS? and calculated the distribution of time span
spent in abortive cycling as a result of these specific rate constants. The
resulted time span was compared to empirical values
\cite{revyakin_abortive_2006} by root of mean square error (RMSE). We repeated
this four steps procedure 100,000 times, and chose the set of parameters that
yielded a minimal RMSE.  

To identify the boundaries of the rate constants, we first performed an
initial screening, marking the rate constants values that have resulted in
large RMSE, in which simulations are inconsistent with experimental data. In
this screening, we found the minimum and maximum values of all rate constants
to be 1 s$^{-1}$ and 25 s$^{-1}$ respectively. 

\begin{figure}
	\begin{center}
        \includegraphics{/home/jorgsk/Dropbox/The-Tome/my_papers/kinetic_initiation/illustrations/model_and_rates.pdf}

	\end{center}
    \caption{Kinetic model scheme and reaction rate constants.  \textbf{A})
    From the open complex (OC) transcription proceeds by nucleotide additions
    cycles from one initial transcribing complex (IC) to the next, where each
    IC is identified in subscript by the length of its nascent RNA. Initial
    transcription proceeds until the nascent RNA has reached the
    experimentally obtained maximum size of abortive transcript (MSAT)
    \cite{hsu_initial_2006}. For ICs with an RNA of 2 nt in length or longer, there is a
    competition between the nucleotide addition cycle (NAC) and backtracking.
    Backstepping leads to a backtracked state from which further backtracking
    and abortive RNA release occurs, returning RNAP to the open complex. From
    the open complex forward transcription may resume once more. \textbf{B})
    Kinetic reactions diagram containing the rate
    constants NAC (denoted here as $k_n$), promoter escape (denoted as $k_e$) and unscrunching and abortive
    release ($k_u$), as well as the rate constants of backstepping, $k_{b,i}$.}
    \label{fig:model_and_rates}
\end{figure}
